\section{Introduction}\label{sec:intro}

Sampling several candidate solutions and letting a learned verifier pick one
(Best-of-$N$) is one of the simplest ways to turn extra inference compute into accuracy
on mathematical reasoning \cite{cobbe2021gsm8k,lightman2023verify,snell2024scaling}. The
verifier is usually a reward model with a few billion parameters. An outcome reward
model (ORM) scores a complete response, and a process reward model (PRM) scores each
reasoning step. In deployment the verifier competes with the generator for accelerator
memory, and an obvious way to make room is post-training quantization (PTQ) of its
weights to 4 bits or fewer \cite{frantar2023gptq,lin2024awq,dettmers2023qlora}.

It is not obvious that this preserves selection quality. PTQ is usually evaluated by
perplexity or by the accuracy of the quantized model on its own task. A verifier, by
contrast, only has to \emph{rank} candidates so that a correct one ends up on top. Small
score perturbations do no harm as long as the winning answer stays the same, but
Best-of-$N$ is sensitive to errors at the top of the ranking
\cite{gao2023overoptimization}. One might therefore expect quantization damage to grow
with the budget $N$ and with problem difficulty, and to matter most for small verifiers.
If, on the other hand, quantized verifiers select as well as full-precision ones, a
fixed memory budget may be better spent on a larger quantized verifier.

We study these questions with frozen candidate pools and analysis plans specified in
advance. Each pool of 32 candidates per problem is generated once, so every verifier and
precision ranks exactly the same responses, and any difference in accuracy comes from
the verifier. We evaluate three ORM sizes and one PRM in BF16, NF4, GPTQ (4 and 3 bits),
AWQ and round-to-nearest on 2,000 GSM8K and MATH problems, and we measure verifier memory
with packed 4-bit kernels. Two follow-up studies, one with a stronger generator and one
with new hard problems, had their plans fixed before their candidates were generated.

Our contributions are as follows.
\begin{enumerate}
\item \textbf{Accuracy equivalence at 4 bits.} GPTQ-W4 meets a prespecified $\pm3$\pp\
equivalence bound for the Best-of-32 accuracy of an 8B ORM (90\% interval $[-0.58,
+1.25]$), and the result holds for both an ORM and a PRM in two follow-up studies.
Three-bit GPTQ with generic calibration degrades selection, most severely for the
smallest Qwen-based verifier and for two Llama-based reward models. In exploratory
tests, calibrating on in-domain candidates reduced the loss from 13.2 to 3.8\pp\ for the
1.7B ORM and from 12.6 to 2.8\pp\ for a Llama-based ORM, while a length-matched generic
control did not help.
\item \textbf{Retained verifier value.} In a follow-up study on 1,060 new hard problems
where both verifiers beat majority vote, a prespecified test shows that 4-bit
quantization keeps more than half of their advantage. The PRM retains an estimated 93\%
of an 8.2\pp\ advantage and the ORM 91\% of a 3.0\pp\ advantage, the latter narrowly
significant. With a strong generator, neither verifier clearly beat majority vote.
\item \textbf{A tested memory comparison.} Under a 6~GiB verifier-only budget and a fixed
workload, a 4B ORM with packed 4-bit weights outperforms the best tested full-precision
verifier that fits. The difference is 2.25\pp\ under the prespecified grader and
1.40\pp\ after a blinded human audit of the winners without an extractable answer.
\item \textbf{Candidate-level instability without accuracy loss.} At 4 bits,
quantization changes the selected candidate on a quarter to a half of problems but
changes the selected answer far less often. Harm and benefit stay small and largely
cancel, which is consistent with answer redundancy in the pools.
\item \textbf{Evaluation lessons.} A strict answer grader can turn reward-model format
preferences into spurious quantization gains. We show how a format/content
decomposition, a lenient grader, a blinded human audit and revised labels detect and correct this.
We also report two prespecified hypotheses that the evaluation data did not support:
that damage concentrates on MATH and that it grows with $N$.
\end{enumerate}
